%! AP Statistics — Unit 9, Lesson 9.4 — Lesson Plan
\documentclass[10pt]{article}
\usepackage{apstats-article}
\usepackage{apstats-boxes}
\usepackage{graphicx}
\graphicspath{{images/}}
\newcommand{\UnitNumberName}{Unit 9: Inference for Quantitative Data: Slopes \quad}
\newcommand{\LessonNumberName}{Lesson 9.4: Setting Up a Test for the Slope of a Regression Model}

\begin{document}
\begin{center}
    {\Large\bfseries \CourseName: \SchoolYear \\
    {\small\bfseries \UnitNumberName \LessonNumberName}}
\end{center}

\begin{tcolorbox}[colback=sky, colframe=navy, boxrule=0.9pt, arc=2mm,
    left=2mm, right=2mm, top=1.5mm, bottom=1.5mm]
    \textbf{Primary Objective:} Students will identify the $t$-test for slope as the
    appropriate procedure for testing a claim about the slope $\beta$ of a population
    regression line, state H$_0$ and H$_a$ in terms of $\beta$, and verify the LINER
    conditions. (Big Idea: VAR)
\end{tcolorbox}

\renewcommand{\arraystretch}{1.6}
\begin{skillbox}[Priority Ideas \& Skills]{goldbox}
\begin{minipage}[t]{0.48\linewidth}
    \textbf{AP Skills Addressed}
    \begin{itemize}
        \item Identify the $t$-test for slope as the appropriate procedure when testing
              a claim about $\beta$ (Skill 1.E; VAR-7.J).
        \item State H$_0$ and H$_a$ in terms of $\beta$ for a $t$-test for slope
              (Skill 1.F; VAR-7.K).
        \item Verify the LINER conditions for inference about slope (Skill 4.C).
    \end{itemize}
\end{minipage}
\hfill
\begin{minipage}[t]{0.48\linewidth}
    \textbf{Key Understandings}
    \begin{itemize}
        \item Hypotheses are written in terms of $\beta$ (population slope), not $b$
              (sample slope).
        \item H$_0$: $\beta = 0$ means no linear relationship exists in the population.
        \item The LINER acronym summarizes five conditions: Linear, Independent, Normal,
              Equal variance, Random.
        \item Residual plots are the primary tool for checking L, N, and E.
        \item A $t$-test for slope parallels the $t$-test for a mean: same distribution
              family, same logic for hypotheses and conditions.
    \end{itemize}
\end{minipage}
\end{skillbox}

\begin{skillbox}[{Vocabulary, Concepts, \& Theorems}]{greenbox}
\renewcommand{\arraystretch}{1.3}
\begin{tabularx}{\linewidth}{lX}
    \textbf{$t$-test for slope} &
        An inference procedure that tests whether the true slope $\beta$ of a population
        regression line equals a specified value $\beta_0$ (usually 0) (VAR-7.J.1) \\
    \textbf{Null hypothesis for slope} &
        H$_0$: $\beta = \beta_0$ (most commonly $\beta_0 = 0$, i.e.\ no linear
        relationship) (VAR-7.K.1) \\
    \textbf{Alternative hypothesis for slope} &
        H$_a$: $\beta \ne \beta_0$, $\beta > \beta_0$, or $\beta < \beta_0$
        (usually $\ne$ on the AP exam) (VAR-7.K.1) \\
    \textbf{$\beta_0$ (hypothesized value)} &
        The value of $\beta$ specified in H$_0$; almost always 0 in AP Statistics \\
    \textbf{LINER conditions} &
        Linear, Independent, Normal, Equal variance, Random --- five conditions
        required for valid inference about the slope of a regression model \\
    \textbf{Residual plot} &
        A scatterplot of residuals vs.\ fitted (or explanatory) values; used to
        check the L, N, and E conditions \\
\end{tabularx}
\end{skillbox}

\begin{skillbox}[Activate Prior Knowledge \& Spiral Review]{sky}
\begin{minipage}[t]{0.48\linewidth}
    \textbf{Prior Knowledge Check (3--4 min)}
    \begin{itemize}
        \item Lesson 9.2: the LINER conditions for a confidence interval for slope
              (same conditions apply to the significance test).
        \item Lesson 9.3: using a CI to make a claim --- if 0 is not captured, there
              is evidence $\beta \ne 0$ --- previews formal test logic.
        \item Unit 7 ($t$-tests for means): H$_0$: $\mu = \mu_0$; H$_a$: $\mu \ne \mu_0$;
              slope replaces mean as the parameter.
    \end{itemize}
\end{minipage}
\hfill
\begin{minipage}[t]{0.48\linewidth}
    \textbf{Connections Forward}
    \begin{itemize}
        \item Computing the $t$-statistic and $p$-value for slope $\to$ Lesson 9.5
        \item Drawing conclusions; full four-step procedure $\to$ Lesson 9.6
        \item Reading computer output for regression inference $\to$ Lesson 9.5
    \end{itemize}
\end{minipage}
\end{skillbox}

\begin{skillbox}[Hook \& Lesson]{sky}
\begin{minipage}[t]{0.48\linewidth}
    \textbf{Hook (5--7 min)}\\
    Display: ``A researcher records elevation (m) and average annual temperature
    ($^\circ$C) for 18 mountain towns and fits a regression line:
    $\hat{y} = 22.4 - 0.0065x$.''\\[4pt]
    Ask: ``She wants to know if elevation \emph{really} predicts temperature in the
    population, or if $b = -0.0065$ could be sampling variability.''\\[4pt]
    Discuss: What three things must she set up before computing anything?
    (Name the procedure, state H$_0$/H$_a$, verify conditions.)
\end{minipage}
\hfill
\begin{minipage}[t]{0.48\linewidth}
    \textbf{Lesson Flow (15 min)}
    \begin{enumerate}
        \item Name the procedure: $t$-test for slope. When is it appropriate?
        \item State H$_0$ and H$_a$ in terms of $\beta$, not $b$.
              H$_0$: $\beta = 0$; H$_a$: $\beta \ne 0$ (most common on AP).
        \item Unpack LINER --- model each letter with the elevation/temperature context.
        \item Practice check: which conditions can be verified from context alone vs.\
              which require the residual plot?
    \end{enumerate}
\end{minipage}
\end{skillbox}

\begin{skillbox}[Explicit Instruction \& Active Monitoring]{sky}
\begin{minipage}[t]{0.48\linewidth}
    \textbf{LINER Checklist}
    \begin{itemize}
        \item \textbf{L (Linear):} Residual plot shows random scatter, no curve.
        \item \textbf{I (Independent):} Random sample or experiment; $n \le 10\%N$
              if sampling without replacement.
        \item \textbf{N (Normal):} Residual histogram approximately bell-shaped,
              or $n > 30$.
        \item \textbf{E (Equal variance):} Residual plot shows roughly constant spread
              across all $x$ values.
        \item \textbf{R (Random):} Data from a random sample or randomized experiment.
    \end{itemize}
\end{minipage}
\hfill
\begin{minipage}[t]{0.48\linewidth}
    \textbf{Common Errors to Address}
    \begin{itemize}
        \item Writing H$_0$: $b = 0$ instead of H$_0$: $\beta = 0$.
        \item Checking linearity from the original scatterplot instead of the
              \emph{residual} plot.
        \item Forgetting to state the 10\% condition for I.
    \end{itemize}
    \textbf{Cold-Call Prompts}
    \begin{itemize}
        \item ``Why do we write $\beta$ in the hypotheses, not $b$?''
        \item ``Which LINER letters require the residual plot? Which can you check
              from context alone?''
    \end{itemize}
\end{minipage}
\end{skillbox}

\begin{skillbox}[Group Work \& Differentiation]{redbox}
\begin{minipage}[t]{0.48\linewidth}
    \textbf{Partner Activity (10--12 min)}\\
    Two scenarios (fiber intake/LDL cholesterol; square footage/sale price). Students
    name the test, define $\beta$ in context, state H$_0$/H$_a$, and verify LINER
    conditions from the given information. Three tiers with increasing autonomy.
\end{minipage}
\hfill
\begin{minipage}[t]{0.48\linewidth}
    \textbf{Differentiation}
    \begin{itemize}
        \item \textit{Support (Tier R)}: Scenario 1 only; sentence frames for
              defining $\beta$; condition checklist with yes/no/N/A choices.
        \item \textit{On-grade (Tier A)}: Both scenarios; write full hypotheses and
              condition justifications.
        \item \textit{Extension (Tier E)}: Add one-sided vs.\ two-sided reasoning
              and argue why two-sided is generally preferred.
        \item \textit{ELL}: Vocabulary card --- ``slope / pendiente,''
              ``residual / residuo,'' ``population / poblaci\'{o}n.''
    \end{itemize}
\end{minipage}
\end{skillbox}

\begin{skillbox}[Individual Work \& Assessment]{redbox}
\begin{minipage}[t]{0.48\linewidth}
    \textbf{Exit Ticket (5 min)}\\
    Context: 35 students, extracurricular hours ($x$) vs.\ GPA ($y$).
    \begin{enumerate}
        \item Name the appropriate inference procedure and justify.
        \item State H$_0$ and H$_a$ in terms of $\beta$; define $\beta$ in context.
        \item Verify LINER conditions given: random scatter in residual plot, equal
              spread, approximately normal residuals, school $N > 350$.
    \end{enumerate}
\end{minipage}
\hfill
\begin{minipage}[t]{0.48\linewidth}
    \textbf{AP-Style Formative Check}\\
    \textit{A researcher fits a regression line to data from a random sample of 25
    adults (age vs.\ resting heart rate) and wants to test whether age is a
    significant linear predictor of heart rate. Which is the correct null hypothesis?}
    \begin{enumerate}[label=(\Alph*)]
        \item H$_0$: $b = 0$
        \item H$_0$: $r = 0$
        \item H$_0$: $\beta = 0$
        \item H$_0$: $\mu = 0$
    \end{enumerate}
    Answer: (C)
\end{minipage}
\end{skillbox}

\begin{skillbox}[Reinforcement \& Extension]{goldbox}
\begin{minipage}[t]{0.48\linewidth}
    \textbf{Homework / Practice}
    \begin{itemize}
        \item Problem 1: Social media hours vs.\ GPA --- name procedure, state
              hypotheses, verify LINER conditions, identify the one condition that
              cannot be fully verified from the given information.
        \item Problem 2: Three contexts requiring H$_0$/H$_a$ and interpretation
              of $\beta = 0$, including one one-sided scenario.
    \end{itemize}
\end{minipage}
\hfill
\begin{minipage}[t]{0.48\linewidth}
    \textbf{Extension / Preview}
    \begin{itemize}
        \item \textit{Extension}: Why is it problematic to choose a one-sided
              alternative after seeing the data? (Discuss data snooping.)
        \item \textit{Preview}: Lesson 9.5 computes the $t$-statistic
              $t = \dfrac{b - \beta_0}{SE_b}$ and finds the $p$-value using
              the $t$-distribution with $df = n - 2$.
    \end{itemize}
\end{minipage}
\end{skillbox}

\end{document}
